\section{Introduction}
\label{sec:intro}

\emph{The question.} Can the Lorentz symmetry of space-time and an internal symmetry (a symmetry
that acts on the components of a field at each point, such as isospin) be joined into one larger
group in which the two kinds of generators do not commute? In four dimensions three theorems answer
\emph{no}. McGlinn~\cite{McGlinn1964} showed that if the internal generators commute with the
translations and form a semisimple Lie algebra, they commute with the Lorentz generators as well.
O'Raifeartaigh~\cite{ORaifeartaigh1965} showed that a finite-dimensional Lie algebra that contains
the Poincar\'e algebra cannot split the masses of a multiplet whose mass is an isolated eigenvalue.
Coleman and Mandula~\cite{ColemanMandula1967} proved the decisive statement: the connected symmetry
group of a nontrivial scattering matrix of particles in four dimensions is locally the direct
product of the Poincar\'e group and a compact internal group. Pelc and Horwitz extended the theorem
to every dimension with one time~\cite{PelcHorwitz1997}. The only known way around it within quantum
field theory with a scattering matrix is to allow anticommuting generators: Haag, {\L}opusza\'nski
and Sohnius~\cite{HLS1975} and Nahm~\cite{Nahm1978} found that supersymmetry is then the unique
extension.

\emph{The task of this paper} is to look for counterexamples to these theorems on the two
eight-dimensional spacetimes of PROJECT spacetimes, the spacetime of signature $7+1$ (seven space
directions and one time) and the spacetime of signature $4+4$, to find the places where their
hypotheses break down in eight dimensions---through the octonions, their non-associativity, their
triality---and to analyse the spin-zero superfluid of the 16-component Dirac field on eight-dimensional
manifolds that combine Lorentz and internal degrees of freedom, instead of on an eight-dimensional
spacetime. We also determine whether this cures the non-renormalizability of the eight-dimensional
theory.

\emph{What a counterexample can be.} A theorem whose hypotheses hold has no counterexample: if it
had one, it would not be a theorem. A counterexample to the \emph{conclusion} ``Lorentz and internal
symmetries cannot be combined'' is therefore a structure in which at least one hypothesis fails, and
the useful content is to say exactly which. That is what we provide. For each hypothesis of the
theorems we construct, on the $7+1$ or the $4+4$ spacetime or on both, an explicit structure that
violates it and in which Lorentz and internal generators are combined, and we verify every statement
about it in exact (rational) arithmetic. Table~\ref{tab:catalogue} in Sec.~\ref{sec:catalogue}
collects them.

\emph{Results.} (i) The octonions and the split octonions, built by the Cayley-Dickson
construction, produce the Dirac matrices of both spacetimes; the $7+1$ metric is the real part of
the square of an octonion and the $4+4$ metric is the norm of a split octonion
(Sec.~\ref{sec:octonions}). (ii) The commutator algebra of the imaginary octonions violates the
Jacobi identity on a set of generators that spans it, so it has no nonzero representation by
operators on a Hilbert space: octonionic charges are excluded by the most basic hypothesis of the
theorems, and their associative envelope is the rotation algebra of eight dimensions itself.
(iii) Cartan's triality~\cite{Cartan1925} exists for the compact form and for the $4+4$ form of the
rotation algebra and exchanges the vector representation with the spinors; for $7+1$ it provably does
not exist (Sec.~\ref{sec:triality}). (iv) On the combined manifolds a four-dimensional spacetime
times a four-dimensional internal fibre the internal group is a part of the eight-dimensional Lorentz
group; on $7+1$ it is compact, on $4+4$ it is non-compact and cannot act unitarily; the automorphism
group $\mathrm G_2$ of the octonions locks the rotations of three-dimensional space to internal
rotations; on $4+4$ a split-octonion multiplication exchanges space-time and internal directions
(Sec.~\ref{sec:combined}). (v) The super-Poincar\'e algebra of $7+1$, with its $U(1)$ R-symmetry
and 128 tensorial charges, satisfies the graded Jacobi identity on every triple of generators and
has positive energy; the one of $4+4$ exists but has no unitary representation with nonzero momentum
(Sec.~\ref{sec:susy}).

(vi) On the combined manifolds a \emph{single} 16-component field forms spin-zero pairs, which is
impossible on the eight-dimensional spacetime: the pairs are vectors and two-forms of the internal
group. Their superfluid has a vacuum threshold, a Cooper instability at finite density, an amplitude
mode at twice the gap, and spontaneous breaking patterns in which an eight-dimensional Lorentz
rotation can be locked to the fermion number (Sec.~\ref{sec:superfluid}). (vii) The combined
manifolds do not cure non-renormalizability: on a compact fibre of radius $R_{\rm NG}$ the tower of
Kaluza-Klein modes restores the eight-dimensional divergence at momenta above $1/R_{\rm NG}$, and on
$4+4$ the tower contains infinitely many tachyons (Sec.~\ref{sec:renorm}).

\emph{Conventions.} The eight coordinates $X=\{x0,\dots,x7\}$ of both spacetimes keep their names;
$x4$ is the time, $x1,x2,x3$ are the three directions of ordinary space, and $x0,x5,x6,x7$ are the
four further directions. Indices $A$, $B$, $C$ run over $0,\dots,7$, the index $j$ over $1,2,3$, the
indices $a$, $b$ over $0,5,6,7$, and $s$, $t$ over $5,6,7$; $K$ denotes a set of indices and $\gamma^K_{\rm NG}$
the ordered product of the matrices $\gamma^A_{\rm NG}$ with $A$ in $K$ (increasing $A$), and $D$ is the number of space
dimensions in which the fields carry momentum (Sec.~\ref{sec:renorm}). Every symbol that this paper
introduces carries the suffix NG, with two kinds of exception that cannot be read as anything else:
the elements of the algebras of Secs.~\ref{sec:octonions}, \ref{sec:triality} and \ref{sec:susy} (octonions,
split octonions, their rotation algebras and the superalgebras) are written in sans-serif type ($\mathsf x$, $\mathsf e_p$,
$\mathsf m_1$, \dots), and $\mathsf 1_8$, $\mathsf 1_{16}$ are unit matrices. The objects of the two
spacetimes keep the tags 1671 ($7+1$) and 1644 ($4+4$) of the project; a symbol with the suffix NG and
without a tag, such as $\gamma^A_{\rm NG}$, denotes the object of the spacetime under discussion
(Sec.~\ref{sec:combined}). Appendix~\ref{app:notation} defines every symbol, Appendix~\ref{app:checks}
lists every machine check, and Appendix~\ref{app:data} the programs and data. Units: $\hbar=1$, the
speed of light 1.
